\chapter{CONSIDERAÇÕES FINAIS}
Este capítulo encerra a sua monografia. Trata-se de uma síntese reflexiva onde o autor fecha o ciclo do trabalho, demonstrando se os objetivos foram alcançados e respondendo definitivamente ao problema de pesquisa.

As considerações finais possuem um raciocínio próprio e devem ser redigidas em texto corrido. Recomenda-se não utilizar citações de outros autores e nem subdividir este capítulo em tópicos ou marcadores.

Para estruturar o seu texto de forma impecável, sugere-se organizar os parágrafos contemplando: síntese dos resultados e resposta ao problema (retomando os resultados mais relevantes e confirmando ou descartando as hipóteses iniciais); contribuições do estudo (benefícios e impactos práticos para a sociedade, instituição investigada e academia); limitações da pesquisa; e propostas para estudos futuros, concluindo com o impacto da realização da monografia na formação enquanto pesquisador(a) e futuro(a) profissional da área.
